% ============================================================
% PRN232 - Section 1: Introduction (Vietnamese)
% ============================================================
\section{\bi{Mở đầu}{Introduction}}
\label{sec:intro}

\bi{
Trong những năm gần đây, nhu cầu theo dõi nội dung truyền hình trực tuyến
tại Việt Nam đã thay đổi rõ rệt. Người dùng không còn chỉ ngồi trước
tivi lúc 21 giờ; họ xem các chương trình theo dõi ngẫu nhiên, bỏ lỡ
lịch phát của kênh mình yêu thích, và muốn được gợi ý nội dung phù hợp
với sở thích. Ở phía bên kia, các đài truyền hình và nền tảng nội
dung số cần một hệ thống lịch phát thống nhất, đáng tin cậy và có thể
mở rộng để phục vụ đồng thời nhiều nền tảng: web, ứng dụng di động và
các thiết bị phát trực tuyến.

Bài báo này trình bày \textbf{OmniCast} --- một hệ thống lập lịch và
phát sóng nội dung số đa nền tảng, được xây dựng cho môn
\textbf{PRN232 --- Building Cross-Platform Back-End Application}.
Hệ thống gồm một back-end NestJS hiệu năng cao phục vụ đồng thời một cổng
web Next.js 15 và một ứng dụng Flutter, với lõi nghiệp vụ tập trung vào
ba năng lực chủ chốt: lịch phát điện tử (EPG) không bao giờ để trống,
cơ chế phát hiện xung đột lịch phát, và trợ lý biên tập AI có kiểm
duyệt theo chuẩn đạo đức nội dung.
}{
In recent years, the way Vietnamese audiences consume television has
changed dramatically. Viewers no longer sit in front of a screen at a
fixed broadcast slot; they browse on-demand, miss the schedules of
channels they follow, and expect content recommendations tailored to
their interests. Broadcasters and digital-content providers, in turn,
need a single, reliable and scalable scheduling system that serves
several delivery surfaces at once: the web, mobile applications and
live-streaming devices.

This paper presents \textbf{OmniCast} --- a multi-platform digital
broadcast scheduling and delivery system built for the course
\textbf{PRN232 --- Building Cross-Platform Back-End Application}. The
system comprises a high-performance NestJS back end that serves both a
Next.js 15 web portal and a Flutter mobile application, with the core
business logic concentrated on three capabilities: an electronic
programme guide (EPG) that is never empty, a broadcast-schedule
conflict-detection mechanism, and an AI editorial assistant with
content-compliance validation.
}

\subsection{\bi{Động cơ và vấn đề cần giải quyết}{Motivation and Problem Statement}}

\bi{
Ba vấn đề kỹ thuật chi phối toàn bộ thiết kế của hệ thống được xác định
từ yêu cầu thực tế của dự án.
}{
Three technical problems dominate the system's design, each derived
from concrete project requirements.
}

\begin{enumerate}[leftmargin=*,itemsep=2pt]
  \item \bi{\textbf{Lịch phát thưa thớt.} Một kênh truyền hình thực tế
  chỉ có khoảng một chương trình gốc mỗi ngày. Nếu dựng lưới EPG 24
  giờ trực tiếp từ bảng dữ liệu, phần lớn lưới sẽ trống và bộ lọc theo
  ngày ở giao diện người dùng trở nên vô nghĩa. Hệ thống cần một cơ
  chế điền lịch tự động, luôn đảm bảo mỗi kênh trong mỗi ngày đều đầy
  đủ 24 giờ nội dung hợp lệ.}{
  \textbf{Sparse schedules.} A real broadcast channel originates only a
  handful of productions per day. Building a 24-hour EPG grid directly
  from the event table leaves most of the grid empty and renders the
  day filter in the user interface meaningless. The system therefore
  requires an automatic filler that guarantees every channel-day grid
  is fully populated with valid content.}

  \item \bi{\textbf{Xung đột lịch phát.} Nếu hai chương trình cùng một
  kênh có khoảng thời gian chồng lấn, người xem sẽ không biết chương
  trình nào được phát. Cơ chế kiểm tra phải phát hiện được cả ba dạng
  giao thoa (đầu, đuôi, bao trùm toàn bộ) và phải chạy trong một câu
  truy vấn cơ sở dữ liệu để tránh điều kiện tranh chấp.}{
  \textbf{Schedule conflicts.} If two programmes on the same channel
  overlap, viewers cannot tell which one is actually on air. The
  validation must detect all three overlap forms (head, tail, full
  containment) and must execute as a single database query to avoid
  race conditions.}

  \item \bi{\textbf{Kiểm duyệt nội dung không tự động.} Quy trình biên
  tập nội dung (mô tả, thẻ, phân loại) là thủ công và tốn thời gian.
  Việc tự động hoá bằng mô hình ngôn ngữ lớn giúp giảm tải, nhưng đòi
  hỏi cơ chế dự phòng khi không có khóa API, và đầu ra phải được kiểm
  tra theo schema trước khi lưu.}{
  \textbf{Manual content moderation.} The editorial workflow (writing
  descriptions, tagging, categorising) is manual and time-consuming.
  Automating it with a large language model reduces the load, but
  requires a fallback path when no API key is present, and the output
  must be schema-validated before it is persisted.}
\end{enumerate}

\subsection{\bi{Đóng góp của bài báo}{Contributions}}

\bi{
Bài báo đóng góp như sau:
\begin{enumerate}[leftmargin=*,itemsep=2pt]
  \item Kiến trúc phân tầng \textbf{NestJS + Prisma + PostgreSQL} với 19
  module nghiệp vụ, 140 endpoint REST và tài liệu OpenAPI tự sinh,
  kèm lớp bảo vệ RBAC bốn vai trò.
  \item Thuật toán phát hiện xung đột lịch phát bằng một truy vấn SQL
  thô, tính toán thời điểm kết thúc hiệu lực bằng
  \texttt{COALESCE(endedAt, scheduledAt + duration)}.
  \item Cơ chế điền lịch EPG 24/7 với lựa chọn ngẫu nhiên có hạt giống
  (seeded PRNG) theo cặp (ngày, kênh), bảo đảm cùng một ngày luôn
  cho cùng một lịch.
  \item Quy trình biên tập AI gọi mô hình một lần duy nhất
  (\texttt{gpt-4o}), kiểm tra schema bằng Zod, có cơ chế thử lại và
  đường dự phòng heuristic.
  \item Một microservice ghi nhật ký kiểm toán giao tiếp qua gRPC,
  tách biệt khỏi tiến trình API chính.
  \item Cổng web Next.js 15 với 33 route và quản trị AI Curator dựa trên
  kiến trúc Next.js App Router.
  \item Quy trình CI/CD năm công việc trên GitHub Actions với cổng chất
  lượng tự động, cùng 482 kiểm thử tự động trên ba bộ kiểm thử.
\end{enumerate}
}{
This paper makes the following contributions.
\begin{enumerate}[leftmargin=*,itemsep=2pt]
  \item A layered \textbf{NestJS + Prisma + PostgreSQL} architecture
  with 19 business modules, 140 REST endpoints and auto-generated
  OpenAPI documentation, protected by a four-role RBAC layer.
  \item A broadcast-conflict detection algorithm implemented as a
  single raw SQL query that computes the effective end time with
  \texttt{COALESCE(endedAt, scheduledAt + duration)}.
  \item A 24/7 EPG filler with seeded pseudo-random selection keyed on
  the (date, channel) pair, guaranteeing that the same day always
  renders the same schedule.
  \item An AI editorial pipeline that issues a single consolidated model
  call (\texttt{gpt-4o}), validates the result against a Zod schema,
  and provides retry plus heuristic-fallback paths.
  \item An audit-logging microservice communicating over gRPC, deployed
  independently of the main API process.
  \item A Next.js 15 web portal with 33 routes and an AI Curator studio
  built on the Next.js App Router architecture.
  \item A five-job GitHub Actions CI/CD pipeline with an automated
  quality gate, covering 482 automated tests across three suites.
\end{enumerate}
}

\subsection{\bi{Phạm vi bài báo}{Scope}}

\bi{
Bài báo tập trung vào phần back-end và cổng web của hệ thống, là hai
sản phẩm của môn PRN232. Phần ứng dụng di động Flutter được trình bày
đầy đủ trong bài báo song song ứng với môn PRM393. Các khía cạnh triển
khai hạ tầng, cấu hình biến môi trường và quy trình phát hành ứng dụng
di động được trình bày trong tài liệu triển khai đi kèm, còn ở đây chỉ
tóm tắt ở mức kiến trúc.}{
This paper concentrates on the back end and the web portal, which are
the two PRN232 deliverables. The Flutter mobile application is
presented separately in the companion PRM393 paper. Infrastructure
deployment, environment-variable configuration and mobile release
procedures are covered in the accompanying deployment documentation
and are only summarised here at the architectural level.
}
